% TOP_PROBLEM: {"problemId": "problem.iwasawa-mu-0-conjecture-for-cyclotomic-z-p-extensions", "canonicalTitle": "Iwasawa mu = 0 Conjecture for Cyclotomic Z_p-Extensions", "releaseRank": 324, "primaryDomain": "number_theory_arithmetic_geometry", "primaryDomainLabel": "Number theory & arithmetic geometry", "displayStatus": "open_with_solved_subcases", "releaseStatus": "open_with_solved_subcases"}
% !TeX program = lualatex
% ATTEMPT_STATUS: unresolved
\documentclass[11pt]{article}
\usepackage[margin=25mm]{geometry}
\usepackage{fontspec}
\setmainfont{Latin Modern Roman}
\usepackage{amsmath,amssymb,mathtools}
\usepackage{unicode-math}
\setmathfont{Latin Modern Math}
\usepackage{xurl,hyperref}
\hypersetup{colorlinks=true,urlcolor=blue}
\setcounter{secnumdepth}{0}
\setlength{\parindent}{0pt}
\setlength{\parskip}{5pt}
\setlength{\emergencystretch}{3em}
\begin{document}
\section*{\texorpdfstring{Iwasawa mu = 0 Conjecture for Cyclotomic Z\_p-Extensions}{Iwasawa mu = 0 Conjecture for Cyclotomic Z\_p-Extensions}}\label{324}
\subsection{Definitions and mathematical statement}
For a number field $F$ and prime $p$, let $F_\infty/F$ be its cyclotomic ${\mathbb{Z}}_p$-extension and $F_n$ the layer of degree $p^n$. Let $A_n$ be the $p$-primary subgroup of the ideal class group of $F_n$. The Iwasawa invariants are characterized, for all sufficiently large $n$, by
\[
 |A_n|=p^{\mu p^n+\lambda n+\nu},
\]
where $\mu,\lambda$ are nonnegative integers and $\nu$ an integer. Equivalently, $\mu$ measures the powers of $p$ in the structure of the torsion module $\varprojlim A_n$ over ${\mathbb{Z}}_p[[\operatorname{Gal}(F_\infty/F)]]\cong{\mathbb{Z}}_p[[T]]$. The conjecture is $\mu=0$ for every $F$ and every $p$. It eliminates the exponential-in-$p^n$ part of the class-group growth; it does not assert $\lambda=0$ or bounded class numbers.
\par\smallskip\textit{Formulation and concept references:} \href{https://arxiv.org/html/2408.07826v2}{[S1]}.

\subsection{Short English statement}
In every cyclotomic prime-adic tower of number fields, is the fastest Iwasawa growth term in the prime part of the class numbers absent?
\subsection{Research attempt: the module obstruction behind cyclotomic class-group growth}
\paragraph{1. What vanishing means and what the source proves.}
The invariant in the question belongs to ideal class groups in the cyclotomic $\mathbb Z_p$-extension of an arbitrary number field. It is not the similarly named invariant of an elliptic-curve Selmer group. Chakravarthy's introduction [S1], checked on 27 September 2026, separates these statements: it gives the classical class-group conjecture, recalls Ferrero--Washington for abelian extensions of $\mathbb Q$, and then turns to elliptic curves. Those abelian results are established external inputs; no reduction of every number field to an abelian field is available here. The present attempt makes the algebraic obstruction explicit and proves several implications and model calculations. It uses the standard structure theorem for finitely generated torsion modules over $\Lambda=\mathbb Z_p[[T]]$ and the usual class-group interpretation of Iwasawa invariants.

\paragraph{2. The elementary factors that carry $\mu$.}
For a finitely generated torsion $\Lambda$-module $X$, the structure theorem gives a pseudo-isomorphism, meaning a homomorphism with finite kernel and cokernel, to an elementary module of the form
\[
 E=\bigoplus_{i=1}^r\Lambda/(p^{a_i})
 \ \oplus\ \bigoplus_{j=1}^s\Lambda/(f_j(T)^{b_j}),
\]
where $a_i,b_j>0$ and the $f_j$ are monic distinguished irreducible polynomials. Distinguished means that every nonleading coefficient is divisible by $p$. The invariants are
\[
 \mu(X)=\sum_i a_i,\qquad
 \lambda(X)=\sum_jb_j\deg f_j.
\]
Thus $\mu=0$ means the absence of the $p$-power elementary factors. It does not require the polynomial factors to disappear.

There is a useful equivalent finiteness test:
\[
 \mu(X)=0\quad\Longleftrightarrow\quad X/pX
 \text{ is finite-dimensional over }\mathbb F_p.
\]
For $E$ this is immediate. A factor $\Lambda/(p^a)$ reduces modulo $p$ to $\mathbb F_p[[T]]$, which has infinite dimension. A factor $\Lambda/(f^b)$ reduces to $\mathbb F_p[T]/(T^{b\deg f})$, which has finite dimension. To transfer the criterion across a pseudo-isomorphism, break the map into two short exact sequences involving its finite kernel and cokernel. Tensoring with $\mathbb F_p$ introduces Tor groups, but for a finite module both its quotient by $p$ and its $p$-torsion are finite. Therefore the resulting kernel and cokernel on reduction modulo $p$ are finite-dimensional. Finite-dimensionality is unchanged.

In the same setting, $\mu(X)=0$ is equivalent to $X$ being finitely generated as a $\mathbb Z_p$-module. If $\mu=0$, Weierstrass division makes each $\Lambda/(f_j^{b_j})$ a finite free $\mathbb Z_p$-module; finite kernel and cokernel preserve finite generation. Conversely, a finitely generated $\mathbb Z_p$-module has finite-dimensional reduction modulo $p$, so the previous criterion applies. This statement is about finite generation, not finite cardinality.

\paragraph{3. Coinvariants exhibit the two different growth mechanisms.}
Let $\omega_n(T)=(1+T)^{p^n}-1$. For $X=\Lambda/(p^a)$,
\[
 X/\omega_nX\cong\Lambda/(p^a,\omega_n).
\]
The polynomial $\omega_n$ is monic of degree $p^n$, and all its nonleading coefficients are divisible by $p$. Weierstrass division makes the quotient a free module of rank $p^n$ over $\mathbb Z/p^a\mathbb Z$. Consequently
\[
 |X/\omega_nX|=p^{a p^n}.
\]
This is exactly the $\mu$ growth scale.

For a different model, choose $c=p$ when $p$ is odd and $c=4$ when $p=2$, and set $Y=\Lambda/(T-c)\cong\mathbb Z_p$. Here $\mu(Y)=0$ and $\lambda(Y)=1$. Its coinvariants have order
\[
 |Y/\omega_nY|=p^{v_p((1+c)^{p^n}-1)}
 =p^{n+v_p(c)}.
\]
The valuation identity can be proved inductively. If $a$ has positive $p$-adic valuation for odd $p$, or valuation at least two for $p=2$, the binomial expansion shows
\[
 v_p((1+a)^p-1)=v_p(a)+1.
\]
The first term $pa$ has that valuation, while every other term has strictly larger valuation under the stated assumptions. Iterate from $a=c$. This proves unbounded linear growth of the exponent while $\mu=0$.

These are algebraic modules, not asserted realizations as class-group inverse limits of specified number fields. They prove why a proposed argument equating $\mu=0$ with bounded class numbers is too strong. Conversely, the factor $\Lambda/(T)$ has infinite coinvariants and is not a valid model for a finite class-group order formula without additional modifications; choosing $T-c$ avoids that degeneracy.

\paragraph{4. What finite differences can detect after stabilization.}
Write $e_n=\log_p|A_n|$. Once the eventual formula is valid at the needed consecutive levels,
\[
 e_{n+1}-e_n=\mu(p-1)p^n+\lambda,
\]
\[
 e_{n+2}-2e_{n+1}+e_n=\mu(p-1)^2p^n.
\]
Thus exact class-group data at three consecutive levels, together with a proved stabilization threshold, can determine $\mu$. If $\mu=0$, the eventual exponent sequence is affine linear. The phrase ``together with'' cannot be deleted: the existence theorem does not supply a universal threshold for every field.

For a purely logical warning, fix any finite prefix of nonnegative integer values and a later index $N$. One can extend that prefix by an eventually affine sequence, or by an eventual sequence $p^n+\lambda n+\nu$, selecting $\nu$ so the late values stay nonnegative. Both sequences satisfy some eventual Iwasawa-shaped formula but have different $\mu$. This does not claim that both arise from number fields. It shows that the eventual-formula theorem alone cannot turn a finite prefix into a universal certificate. Arithmetic information controlling the onset or the inverse-limit module is needed.

\paragraph{5. Passing to a cyclotomic layer does not remove a positive $\mu$.}
Take the $m$th layer $F_m$ as a new base field, and index its cyclotomic extension from zero. Its $n$th layer is $F_{m+n}$. Hence the original eventual formula gives
\[
 e_{m+n}=\mu p^m p^n+\lambda n+(\lambda m+\nu).
\]
Uniqueness of the exponential, linear, and constant coefficients implies
\[
 \mu(F_m)=p^m\mu(F),\qquad \lambda(F_m)=\lambda(F).
\]
Uniqueness follows by subtracting two formulas and dividing by $p^n$, then by $n$: the successive coefficients must vanish. In particular $\mu(F_m)=0$ if and only if $\mu(F)=0$. Moving sufficiently far up the same tower cannot kill a positive $\mu$; it increases its coefficient under the new indexing. This is an important normalization check when comparing computations made with different base layers.

\paragraph{6. A genuine descent statement for prime-to-$p$ extensions.}
Let $E/F$ be a finite extension of degree $d$ prime to $p$. Its intersection with the cyclotomic $\mathbb Z_p$-extension of $F$ is $F$, because the intersection degree divides both $d$ and a power of $p$. The cyclotomic layers of $E$ are $E_n=EF_n$, and $[E_n:F_n]=d$. Extension of ideals and norm induce maps on the $p$-primary class groups satisfying
\[
 A(F_n)\xrightarrow{j}A(E_n)\xrightarrow{N}A(F_n),
 \qquad N\circ j=d.
\]
Multiplication by $d$ is an automorphism of every finite $p$-group, so $j$ is injective. Therefore $e_n(F)\leq e_n(E)$. Dividing the eventual growth formulas by $p^n$ and taking limits gives $\mu(F)\leq\mu(E)$. In particular vanishing over $E$ implies vanishing over $F$.

This is a useful descent rule, but it does not place an arbitrary nonabelian field inside an abelian extension of $\mathbb Q$: every subfield of an abelian Galois extension is itself abelian over $\mathbb Q$. Nor does the proof apply unchanged to extensions of degree divisible by $p$, where multiplication by the degree need not be injective on class groups. Both restrictions matter when attempting to use Ferrero--Washington as a universal starting point.

\paragraph{7. The actual arithmetic gap.}
Let $X$ be the standard inverse limit of the $p$-primary class groups under norm maps. The structure theorem and the class-group growth theorem identify the relevant $\mu$ with that of $X$. The reduction in Section 2 says that it would suffice to prove finite-dimensionality of $X/pX$, or finite generation of $X$ over $\mathbb Z_p$, for every cyclotomic tower. It does not prove either one. Norm relations, ramification information, and finite-level class groups must still rule out a factor behaving like $\Lambda/(p)$.

Replacing class groups by another Iwasawa module with a similar notation does not bridge this gap. A main conjecture relating a characteristic ideal to a $p$-adic $L$-function can help only when it applies to this module and the required $\mu$-vanishing of that analytic object is also established.

\paragraph{8. Diagnostics and outcome.}
The saved exact calculations verify the binomial valuation identity in the displayed odd-prime and $p=2$ models, the finite-difference formulas, and the change-of-base-layer coefficients. These are module-model and indexing checks, not computed class groups for a range of fields.

\textbf{UNRESOLVED.} The attempt gives a precise module criterion for $\mu=0$, distinguishes it from $\lambda=0$, proves prime-to-$p$ descent, and rules out an invalid tower-shift shortcut. The first general gap remains the arithmetic exclusion of $p$-power elementary factors in the cyclotomic class-group module of an arbitrary number field.

\subsection{Sources}
\begin{itemize}
\item[S1]Adithya Chakravarthy, The Iwasawa mu-invariants of Elliptic Curves over Q, arXiv:2408.07826v2 (2024), Introduction, Conjecture 1.1 and following paragraph.
\url{https://arxiv.org/html/2408.07826v2}.
\textit{Formulation source.}
\end{itemize}
\end{document}
